\documentclass[11pt, a4paper]{article}

% Required Packages
\usepackage{amsmath, amsthm, amssymb, amsfonts}
\usepackage{geometry}
\usepackage{graphicx}
\usepackage{hyperref}
\usepackage{natbib}
\usepackage{bm} % For bold math vectors
\usepackage{xcolor}

\geometry{margin=1in}

% Custom Commands
\newcommand{\pderiv}[2]{\frac{\partial #1}{\partial #2}}
\newcommand{\curl}{\nabla \times}
\renewcommand{\div}{\nabla \cdot}

\title{\textbf{Mathematical Foundations of the HopfNet-MHD Engine}}
\author{Mohana Rangan Desigan}
\date{\today}

\begin{document}

\maketitle

\begin{abstract}
This document outlines the rigorous mathematical formulation governing the 3D Pseudo-Spectral Hall-MHD engine. It derives the uncurled induction equations, the Coulomb gauge projection in Fourier space, the regularized analytical construction of a magnetic Hopf link, the ETDRK4 contour integration scheme, and the conservation of Generalized Helicity.
\end{abstract}

\tableofcontents
\newpage

\section{The Governing Equations: Hall-MHD in Vector Potential}
\label{sec:equations}

To computationally simulate a plasma, we must translate the continuous laws of electromagnetism into discrete mathematical steps. The most intuitive starting point is Faraday's Law of Induction, which dictates how magnetic fields ($\mathbf{B}$) evolve over time in response to an electric field ($\mathbf{E}$):
\begin{equation}
    \frac{\partial \mathbf{B}}{\partial t} = -\nabla \times \mathbf{E}
\end{equation}
However, using Equation (1) as the primary engine for our simulation introduces a fatal algorithmic vulnerability related to Gauss's Law for Magnetism. Maxwell's equations mandate that magnetic field lines have no beginning or end; there are no magnetic monopoles. Mathematically, this is expressed as the solenoidal constraint: $\nabla \cdot \mathbf{B} = 0$. 

In a discrete computer simulation, floating-point truncation errors inevitably accumulate at each time-step. If we evolve $\mathbf{B}$ directly, these tiny errors cause the divergence of the magnetic field to drift away from zero over time ($\nabla \cdot \mathbf{B} \neq 0$). Physically, this implies the computer has artificially created magnetic monopoles, which exert fictitious forces parallel to the field lines and catastrophically destroy the topological validity of the plasma's geometry.

\subsection{The Vector Potential Safeguard}
To natively preserve the $\nabla \cdot \mathbf{B} = 0$ constraint to absolute machine precision, we must reformulate the dynamics. Instead of tracking the magnetic field, we track the \textit{magnetic vector potential}, a field denoted as $\mathbf{A}$, which is defined mathematically such that its curl yields the magnetic field:
\begin{equation}
    \mathbf{B} = \nabla \times \mathbf{A}
\end{equation}
By standard vector calculus identities, the divergence of the curl of any vector field is identically zero ($\nabla \cdot (\nabla \times \mathbf{A}) \equiv 0$). Therefore, by simulating the evolution of $\mathbf{A}$ instead of $\mathbf{B}$, we guarantee that no magnetic monopoles can ever be computationally generated, effectively hardcoding Gauss's Law into the algorithmic architecture.

\subsection{The Generalized Ohm's Law in Plasma}
To determine how $\mathbf{A}$ evolves, we must define the electric field $\mathbf{E}$ inside the plasma. In IB Physics, Ohm's Law is taught as $V=IR$, or in continuous form, $\mathbf{J} = \sigma \mathbf{E}$, where $\mathbf{J}$ is the current density and $\sigma$ is conductivity. However, a plasma is not a static copper wire; it is a violently moving fluid where moving magnetic fields generate their own internal forces. 

In dimensionless units (where vacuum permeability $\mu_0 = 1$), the electric field generated within a two-fluid Hall-MHD plasma is described by the Generalized Ohm's Law:
\begin{equation}
    \mathbf{E} = -\underbrace{\mathbf{v} \times \mathbf{B}}_{\text{Ideal Advection}} + \underbrace{\eta \mathbf{J}}_{\text{Resistive Drag}} + \underbrace{\frac{d_i}{\rho}(\mathbf{J} \times \mathbf{B})}_{\text{The Hall Term}} - \underbrace{\frac{d_i}{\rho} \nabla P_e}_{\text{Electron Pressure}}
\end{equation}
Each term in Equation (3) represents a distinct physical causality driving the plasma:
\begin{itemize}
    \item \textbf{Ideal Advection ($-\mathbf{v} \times \mathbf{B}$):} This is the standard Lorentz force. It locks the magnetic field lines to the bulk fluid velocity ($\mathbf{v}$), dragging them together.
    \item \textbf{Resistive Drag ($\eta \mathbf{J}$):} Here, $\eta$ is the fluid resistivity and $\mathbf{J} = \nabla \times \mathbf{B}$ is the current density. This accounts for electron-ion collisions. It is the term that breaks the "frozen-in" constraint, allowing magnetic fields to diffuse and tear.
    \item \textbf{The Hall Term ($\frac{d_i}{\rho}(\mathbf{J} \times \mathbf{B})$):} This term elevates the simulation from standard MHD to Hall-MHD. It is parameterized by the density ($\rho$) and the \textit{ion inertial length} ($d_i$)—the physical scale at which heavy ions become too sluggish to follow light, fast electrons. Because the species decouple at this microscopic scale, the electrons carry the magnetic field lines away much faster than the bulk fluid, generating a secondary electric field. This is the mechanism responsible for the ultra-fast disruption times observed in real Tokamaks.
    \item \textbf{Electron Pressure ($-\frac{d_i}{\rho} \nabla P_e$):} The spatial gradient of electron thermal pressure ($P_e$), which pushes charges outward, creating an electrostatic field.
\end{itemize}

\subsection{Uncurling the Induction Equation}
To find the governing equation for the vector potential, we substitute the Generalized Ohm's Law (Eq. 3) into Faraday's Law (Eq. 1), and simultaneously substitute $\mathbf{B} = \nabla \times \mathbf{A}$ into the time derivative:
\begin{equation}
    \frac{\partial (\nabla \times \mathbf{A})}{\partial t} = -\nabla \times \left[ -\mathbf{v} \times \mathbf{B} + \eta \mathbf{J} + \frac{d_i}{\rho}(\mathbf{J} \times \mathbf{B}) - \frac{d_i}{\rho} \nabla P_e \right]
\end{equation}

Because partial differentiation with respect to time commutes with spatial differentiation (the curl operator), we can pull the $\nabla \times$ operator outside the time derivative on the left-hand side. Distributing the negative sign on the right-hand side yields:
\begin{equation}
    \nabla \times \left( \frac{\partial \mathbf{A}}{\partial t} \right) = \nabla \times \left[ \mathbf{v} \times \mathbf{B} - \eta \mathbf{J} - \frac{d_i}{\rho}(\mathbf{J} \times \mathbf{B}) + \frac{d_i}{\rho} \nabla P_e \right]
\end{equation}

We now have an equation of the form $\nabla \times \mathbf{X} = \nabla \times \mathbf{Y}$. Moving all terms to the left side isolates the curl operator:
\begin{equation}
    \nabla \times \left[ \frac{\partial \mathbf{A}}{\partial t} - \mathbf{v} \times \mathbf{B} + \eta \mathbf{J} + \frac{d_i}{\rho}(\mathbf{J} \times \mathbf{B}) - \frac{d_i}{\rho} \nabla P_e \right] = 0
\end{equation}

In basic IB Physics, we learn that if the curl of an electric field is zero, the field is conservative and can be perfectly described as the gradient of a scalar voltage ($- \nabla V$). The same mathematical principle—formally known as Helmholtz's Theorem—applies here. Because the curl of the bracketed expression is identically zero, the entire expression must equal the gradient of some arbitrary scalar potential field, which we will define as $-\nabla \phi$. 

Thus, we can "uncurl" the equation by dropping the $\nabla \times$ operator and equating the internal terms to this scalar potential:
\begin{equation}
    \frac{\partial \mathbf{A}}{\partial t} - \mathbf{v} \times \mathbf{B} + \eta \mathbf{J} + \frac{d_i}{\rho}(\mathbf{J} \times \mathbf{B}) - \frac{d_i}{\rho} \nabla P_e = -\nabla \phi
\end{equation}

\subsection{The Barotropic Assumption and Final Form}
Equation (7) is almost complete, but the electron pressure gradient term is computationally inconvenient. To simplify it, we make a standard physical assumption: the plasma is \textit{barotropic}. This means the electron pressure depends solely on the density, $P_e = P_e(\rho)$. Under this assumption, thermodynamics dictates that the term $\frac{1}{\rho}\nabla P_e$ acts as the exact gradient of the fluid's specific enthalpy, $\nabla h_e$. 

Because the gradient of a scalar (enthalpy) added to the gradient of another scalar (our arbitrary $\phi$) simply results in a new composite scalar gradient, the messy electron pressure term is mathematically absorbed entirely into our generalized scalar potential $-\nabla \phi$. 

Rearranging the terms to isolate the time evolution of the vector potential, we arrive at the final induction equation that our C++ physics engine will solve:
\begin{equation}
    \frac{\partial \mathbf{A}}{\partial t} = \mathbf{v} \times \mathbf{B} - \frac{d_i}{\rho} (\mathbf{J} \times \mathbf{B}) - \eta \mathbf{J} - \nabla \phi
    \label{eq:hall_mhd_A}
\end{equation}

Equation (\ref{eq:hall_mhd_A}) elegantly and causally defines the magnetic dynamics of the system. However, a critical ambiguity remains: the scalar potential $\phi$ is undetermined. Because $\mathbf{A}$ is a mathematical construct, there are infinitely many valid vector potentials that yield the same physical magnetic field $\mathbf{B}$—a concept known as gauge freedom. To solve the equation deterministically on a computer, we must choose a specific gauge constraint. The selection of $\phi$ and the mathematical projection of this equation into the Coulomb Gauge ($\nabla \cdot \mathbf{A} = 0$) will be addressed in Section \ref{sec:gauge}.

\section{The Coulomb Gauge and Spectral Projection}
\label{sec:gauge}

As established in Equation (\ref{eq:hall_mhd_A}), the evolution of the vector potential $\mathbf{A}$ is governed by the physical forces within the plasma, but it remains mathematically underdetermined due to the arbitrary scalar potential $\phi$. To resolve this ambiguity and close the system of equations, we must impose a gauge condition. 

For incompressible and weakly compressible magneto-fluid dynamics, the most robust choice is the \textbf{Coulomb Gauge}, defined mathematically as:
\begin{equation}
    \nabla \cdot \mathbf{A} = 0
\end{equation}
Physically, this ensures that the vector potential is purely solenoidal, containing no longitudinal (compressive) components. However, enforcing this constraint in a standard finite-difference computer simulation usually requires solving a computationally expensive Poisson equation for $\phi$ at every single time-step, which creates a massive algorithmic bottleneck. 

To bypass this, we leverage the power of the \textit{pseudo-spectral} method. By translating our physical universe into Fourier space, differential calculus is transformed into algebraic geometry, allowing us to enforce the Coulomb gauge perfectly and instantaneously.

\subsection{Transformation into the Spectral Domain}
We represent the vector potential $\mathbf{A}(\mathbf{x}, t)$ as a sum of continuous spatial frequencies using the discrete Fourier transform:
\begin{equation}
    \mathbf{A}(\mathbf{x}, t) = \sum_{\mathbf{k}} \hat{\mathbf{A}}(\mathbf{k}, t) e^{i \mathbf{k} \cdot \mathbf{x}}
\end{equation}
where $\hat{\mathbf{A}}$ is the complex Fourier coefficient and $\mathbf{k} = (k_x, k_y, k_z)$ is the wavevector. 

In the spectral domain, the spatial gradient operator ($\nabla$) corresponds to multiplication by $i\mathbf{k}$. Therefore, the Coulomb gauge constraint ($\nabla \cdot \mathbf{A} = 0$) translates elegantly into a geometric dot product in $k$-space:
\begin{equation}
    i\mathbf{k} \cdot \hat{\mathbf{A}} = 0 \quad \implies \quad \mathbf{k} \cdot \hat{\mathbf{A}} = 0
\end{equation}
Equation (11) dictates that for the Coulomb gauge to hold, the Fourier coefficients of the vector potential must be strictly orthogonal to their corresponding wavevectors. The field is purely transverse.

\subsection{Absorbing the Scalar Potential}
To see how this geometric constraint eliminates the need to calculate $\phi$, let us group all the explicitly known physical terms (Advection, Hall, and Resistive drag) from Equation (\ref{eq:hall_mhd_A}) into a single "raw" forcing vector $\mathbf{R}$:
\begin{equation}
    \mathbf{R} = \mathbf{v} \times \mathbf{B} - \frac{d_i}{\rho} (\mathbf{J} \times \mathbf{B}) - \eta \mathbf{J}
\end{equation}
Equation (\ref{eq:hall_mhd_A}) can now be written compactly as:
\begin{equation}
    \frac{\partial \mathbf{A}}{\partial t} = \mathbf{R} - \nabla \phi
\end{equation}

Taking the Fourier transform of this entire equation yields the temporal evolution in $k$-space:
\begin{equation}
    \frac{\partial \hat{\mathbf{A}}}{\partial t} = \hat{\mathbf{R}} - i\mathbf{k}\hat{\phi}
\end{equation}

To enforce the Coulomb gauge for all time, we require that the time derivative of the divergence of $\mathbf{A}$ is also zero ($\frac{\partial}{\partial t}(\nabla \cdot \mathbf{A}) = 0$). In Fourier space, this requires $\mathbf{k} \cdot \frac{\partial \hat{\mathbf{A}}}{\partial t} = 0$. 

Taking the dot product of the wavevector $\mathbf{k}$ with Equation (14), we obtain:
\begin{equation}
    \mathbf{k} \cdot \frac{\partial \hat{\mathbf{A}}}{\partial t} = \mathbf{k} \cdot \hat{\mathbf{R}} - \mathbf{k} \cdot (i\mathbf{k}\hat{\phi}) = 0
\end{equation}
Because $\mathbf{k} \cdot \mathbf{k} = |\mathbf{k}|^2$, we can solve directly for the unknown scalar potential term $i\hat{\phi}$:
\begin{equation}
    i\hat{\phi} = \frac{\mathbf{k} \cdot \hat{\mathbf{R}}}{|\mathbf{k}|^2}
\end{equation}

\subsection{The Orthogonal Projection Tensor}
We have just algebraically isolated the exact contribution of the scalar potential. Substituting Equation (16) back into Equation (14), we get:
\begin{equation}
    \frac{\partial \hat{\mathbf{A}}}{\partial t} = \hat{\mathbf{R}} - \mathbf{k} \left( \frac{\mathbf{k} \cdot \hat{\mathbf{R}}}{|\mathbf{k}|^2} \right)
\end{equation}

By factoring out the raw physical forcing vector $\hat{\mathbf{R}}$, we can rewrite this operation using the identity matrix $\mathbb{I}$ and the outer product of the wavevectors:
\begin{equation}
    \frac{\partial \hat{\mathbf{A}}}{\partial t} = \underbrace{\left( \mathbb{I} - \frac{\mathbf{k} \otimes \mathbf{k}}{|\mathbf{k}|^2} \right)}_{\mathcal{P}(\mathbf{k})} \hat{\mathbf{R}}
\end{equation}

This reveals the fundamental architectural elegance of the spectral method. The operator acting on $\hat{\mathbf{R}}$ is the \textbf{Orthogonal Projection Tensor}, denoted $\mathcal{P}(\mathbf{k})$. In index notation, it is written as:
\begin{equation}
    \mathcal{P}_{ij}(\mathbf{k}) = \delta_{ij} - \frac{k_i k_j}{|\mathbf{k}|^2}
\end{equation}

\subsection{Causal Conclusion}
Equation (18) is the computational crown jewel of our Phase 1 engine. It mathematically proves that we do \textit{not} need to evaluate the scalar potential $\phi$ or solve a Poisson equation. 

By simply calculating the raw physical forces ($\mathbf{R}$) in real space, transforming them into Fourier space ($\hat{\mathbf{R}}$), and multiplying them by the projection tensor $\mathcal{P}_{ij}$, we perfectly strip away any longitudinal components. The resulting vector $\frac{\partial \hat{\mathbf{A}}}{\partial t}$ is guaranteed to be strictly orthogonal to $\mathbf{k}$. When stepped forward in time, this mathematically mandates that $\nabla \cdot \mathbf{A} = 0$ is preserved to absolute floating-point precision ($10^{-15}$), ensuring that no unphysical magnetic monopoles can ever corrupt the topological decay of the Hopf Link.

\section{Initial Conditions: The 3D Magnetic Hopf Link}
\label{sec:hopf_link}

With the mathematical machinery to strictly enforce the $\nabla \cdot \mathbf{A} = 0$ constraint established, we must now define the initial state of the simulation ($t=0$). The overarching physical objective of this research is to predict topological collapse. Therefore, we must initialize the plasma in a highly stressed topological state that is fundamentally compelled to decay.

We choose the \textbf{Magnetic Hopf Link}: two distinct, interlocking magnetic flux rings.

\subsection{Physical Motivation: Helicity and Taylor Relaxation}
In plasma physics, the topological complexity of a magnetic field—its knottedness, twisting, and linking—is quantified by the Magnetic Helicity integral, $H = \int \mathbf{A} \cdot \mathbf{B} \, dV$. 

According to Taylor Relaxation theory, a turbulent plasma will spontaneously seek the lowest possible energy state (a force-free state where $\mathbf{J} \times \mathbf{B} = 0$). However, it must do so while conserving its global magnetic helicity. A magnetic Hopf link carries a global linking number of $L_k = 1$, meaning $H \neq 0$. Because the two rings are interlocked, they represent a state of high magnetic potential energy. The plasma \textit{wants} to separate the rings to minimize energy, but to do so, the magnetic field lines must tear and pass through one another—a process governed strictly by magnetic reconnection. 

By initializing a Hopf link, we causally guarantee the formation of a localized reconnection singularity (the current sheet) as the system attempts to relax.

\subsection{The Biot-Savart Formulation and the Singularity Problem}
To computationally initialize this configuration, we must define the vector potential $\mathbf{A}(x,y,z)$ at every point in our 3D grid. We construct this using the Biot-Savart Law. For a steady current $I$ traveling along a closed 1D contour $C$, the vector potential at a position $\mathbf{r}$ is:
\begin{equation}
    \mathbf{A}(\mathbf{r}) = \frac{\mu_0 I}{4\pi} \oint_C \frac{d\mathbf{l}'}{|\mathbf{r} - \mathbf{r}'|}
\end{equation}

\textbf{The Computational Vulnerability:} Equation (20) presents a catastrophic problem for a pseudo-spectral engine. An infinitely thin current wire causes the distance denominator $|\mathbf{r} - \mathbf{r}'|$ to approach zero at the core, creating a mathematical singularity ($\mathbf{A} \to \infty$). In Fourier space, a singularity (or a step-function) requires infinite frequencies to resolve. Because our $128^3$ grid has a strict Nyquist frequency limit, injecting a singularity will cause severe \textit{Gibbs ringing}—unphysical, high-frequency oscillations that will instantly trigger numerical instability and destroy the hyper-diffusion term ($-\nu k^4$).

To bypass this computational hazard, we apply a mathematical regularization. We replace the infinitely thin wire with a flux tube of finite core radius $a$. This acts as a spatial low-pass filter, smoothing the current distribution into a Gaussian profile and guaranteeing that the spectral energy drops to zero well before the grid's Nyquist limit. The desingularized Biot-Savart integral becomes:
\begin{equation}
    \mathbf{A}_{reg}(\mathbf{r}) = \frac{\mu_0 I}{4\pi} \oint_C \frac{d\mathbf{l}'}{\sqrt{|\mathbf{r} - \mathbf{r}'|^2 + a^2}}
\end{equation}

\subsection{Explicit Parametrization of the Hopf Link}
We define two orthogonal current loops, $C_1$ and $C_2$, with equal radii $R$, carrying current $I_0$. To ensure they interlock symmetrically within the center of our triply periodic domain $[-L, L]^3$, we offset them by a distance $d$ along the $y$-axis (where $0 < d < R$).

\textbf{Ring 1 ($C_1$)} is placed in the $xy$-plane, centered at $(0, d, 0)$. It is parametrized by the angle $\theta \in [0, 2\pi)$:
\begin{equation}
    \mathbf{r}_1(\theta) = \langle R\cos\theta, \; d + R\sin\theta, \; 0 \rangle, \quad d\mathbf{l}_1 = \langle -R\sin\theta, \; R\cos\theta, \; 0 \rangle \, d\theta
\end{equation}

\textbf{Ring 2 ($C_2$)} is placed in the $yz$-plane, centered at $(0, -d, 0)$. It is parametrized by the angle $\phi \in [0, 2\pi)$:
\begin{equation}
    \mathbf{r}_2(\phi) = \langle 0, \; -d + R\cos\phi, \; R\sin\phi \rangle, \quad d\mathbf{l}_2 = \langle 0, \; -R\sin\phi, \; R\cos\phi \rangle \, d\phi
\end{equation}

By substituting the parameterizations from Equations (22) and (23) into the regularized Biot-Savart Law (Eq. 21), we derive the explicit, continuous integrals for the initial vector potential $\mathbf{A}_{total} = \mathbf{A}_1 + \mathbf{A}_2$. We define the regularized distance functions:
\begin{align}
    D_1(\mathbf{r}, \theta) &= \sqrt{(x - R\cos\theta)^2 + (y - d - R\sin\theta)^2 + z^2 + a^2} \\
    D_2(\mathbf{r}, \phi)   &= \sqrt{x^2 + (y + d - R\cos\phi)^2 + (z - R\sin\phi)^2 + a^2}
\end{align}

The three spatial components of the magnetic vector potential evaluated at any grid coordinate $(x,y,z)$ are:
\begin{align}
    A_x(x,y,z) &= \frac{\mu_0 I_0}{4\pi} \int_0^{2\pi} \frac{-R\sin\theta}{D_1(\mathbf{r}, \theta)} \, d\theta \\
    A_y(x,y,z) &= \frac{\mu_0 I_0}{4\pi} \left[ \int_0^{2\pi} \frac{R\cos\theta}{D_1(\mathbf{r}, \theta)} \, d\theta + \int_0^{2\pi} \frac{-R\sin\phi}{D_2(\mathbf{r}, \phi)} \, d\phi \right] \\
    A_z(x,y,z) &= \frac{\mu_0 I_0}{4\pi} \int_0^{2\pi} \frac{R\cos\phi}{D_2(\mathbf{r}, \phi)} \, d\phi
\end{align}

\subsection{Algorithmic Initialization}
During the initialization phase (Step 0) of the C++ physics engine, these three 1D integrals are evaluated at every point on the $128^3$ grid using Simpson's Rule or Gaussian quadrature. 

Because the superposition of any two valid vector potentials remains a valid vector potential, this construction inherently models the true magnetic field without introducing fictitious divergence. To ensure absolute spectral purity before the physics loop begins, the resulting discrete tensor $\mathbf{A}(x,y,z)$ is immediately passed through the Fourier projection tensor $\mathcal{P}_{ij}(\mathbf{k})$ derived in Section \ref{sec:gauge}, guaranteeing the Coulomb gauge constraint is flawlessly locked in.

\section{Exponential Time Differencing (ETDRK4) and Stiffness}
\label{sec:etdrk4}

To accurately resolve the microscopic topological tearing of the Hopf link without injecting unphysical grid-scale noise, we must apply a \textit{hyper-resistivity} and \textit{hyper-viscosity} term to our equations. Instead of standard Laplacian diffusion ($\nabla^2$), we apply a bi-harmonic operator ($-\nu_4 \nabla^4$). In the spectral domain, this operator scales as $-k^4$. 

Physically, a $-k^4$ operator acts as a "spectral brick wall." At low wavenumbers (the macroscopic plasma bulk), $k^4 \approx 0$, ensuring the fluid behaves as an ideal, collisionless plasma. At high wavenumbers (approaching the grid's Nyquist limit), $k^4 \to \infty$, aggressively dissipating cascaded turbulent energy before it can alias and violate thermodynamics.

\subsection{The CFL Violation and the Stiffness Problem}
While physically ideal, this hyper-dissipation introduces extreme mathematical \textit{stiffness}. In any explicit time-stepping scheme (like standard Runge-Kutta 4), the maximum stable timestep $\Delta t$ is strictly limited by the fastest timescale in the system (the Courant-Friedrichs-Lewy or CFL condition). Because our dissipation scales as $k^4$, the required timestep for a $128^3$ grid would scale as $\Delta x^4 / \nu_4$, pushing $\Delta t \sim 10^{-8}$. Integrating a 5-second plasma disruption at this timestep would require billions of iterations, rendering the simulation computationally impossible for a 16GB desktop architecture.

To bypass the CFL limit, we employ the \textbf{Exponential Time Differencing (ETD)} method. We begin by splitting our governing spectral ODE (Equation 18) into a linear (stiff) operator $\mathcal{L}$ and a nonlinear (non-stiff) operator $\mathcal{N}$:
\begin{equation}
    \frac{\partial \hat{\mathbf{A}}}{\partial t} = \mathcal{L}\hat{\mathbf{A}} + \mathcal{N}(\hat{\mathbf{A}}, t)
\end{equation}
where $\mathcal{L} = -\eta k^2 - \nu_4 k^4$, and $\mathcal{N}$ contains the advection and Hall terms. 

Instead of approximating the stiff linear derivative, ETD multiplies the entire equation by the integrating factor $e^{-\mathcal{L}t}$ and integrates exactly over a timestep $\Delta t$:
\begin{equation}
    \hat{\mathbf{A}}(t + \Delta t) = e^{\mathcal{L}\Delta t}\hat{\mathbf{A}}(t) + e^{\mathcal{L}\Delta t} \int_0^{\Delta t} e^{-\mathcal{L}\tau} \mathcal{N}(\hat{\mathbf{A}}(t+\tau), t+\tau) \, d\tau
\end{equation}
This formulation is profoundly powerful: the linear stiffness is solved analytically (which is unconditionally stable for any $\Delta t$), leaving only the non-stiff nonlinear integral to be approximated via Runge-Kutta.

\subsection{The Floating-Point Singularity ($0/0$)}
To evaluate the nonlinear integral numerically using a 4th-order Runge-Kutta scheme (ETDRK4), we must calculate specific coefficient matrices. The simplest of these, used for the first-order Euler step, is:
\begin{equation}
    f_1(\mathcal{L}) = \frac{e^{\mathcal{L}\Delta t} - 1}{\mathcal{L}}
\end{equation}

\textbf{The Computational Vulnerability:} Equation (31) presents a catastrophic floating-point vulnerability. For the zero-mode ($k=0$), which represents the mean background magnetic field, the linear operator $\mathcal{L} = 0$. Consequently, Equation (31) evaluates strictly to $0/0$. Even for small wavenumbers ($k \approx 1$), $\mathcal{L}$ is very close to zero, causing severe \textit{catastrophic cancellation} in double-precision arithmetic. A single `NaN` (Not-a-Number) generated at $k=0$ will instantly infect the entire $128^3$ grid upon the next inverse Fourier transform, destroying the simulation.

\subsection{The Kassam-Trefethen Contour Integration}
To cure this $0/0$ singularity without resorting to clumsy Taylor series expansions (which suffer from their own truncation errors), we employ Cauchy's Integral Formula from complex analysis, famously applied to ETDRK4 by Kassam and Trefethen (2005).

Cauchy's theorem states that any analytic function of an operator $f(\mathcal{L})$ can be evaluated by integrating over a closed contour $\Gamma$ in the complex plane that encloses the eigenvalues of $\mathcal{L}$:
\begin{equation}
    f(\mathcal{L}) = \frac{1}{2\pi i} \oint_\Gamma \frac{f(z)}{z\mathbb{I} - \mathcal{L}} \, dz
\end{equation}
By moving the operator $\mathcal{L}$ into the denominator of the resolvent $(z - \mathcal{L})^{-1}$, the $0/0$ cancellation is mathematically annihilated. The function is evaluated on a contour safely away from the origin.

We parameterize the contour $\Gamma$ as a circle of radius $R$ centered at the origin: $z = R e^{i\theta}$, implying $dz = i R e^{i\theta} d\theta$. Substituting this into Cauchy's formula yields:
\begin{equation}
    f(\mathcal{L}) = \frac{1}{2\pi i} \int_0^{2\pi} \frac{f(R e^{i\theta})}{R e^{i\theta} - \mathcal{L}} i R e^{i\theta} \, d\theta = \ \frac{1}{2\pi} \int_0^{2\pi} \frac{f(R e^{i\theta})}{R e^{i\theta} - \mathcal{L}} R e^{i\theta} \, d\theta
\end{equation}

\subsection{Discrete Riemann Sum Implementation}
For implementation in our C++ engine, we discretize the continuous integral into a finite Riemann sum over $M$ equally spaced points along the circular contour. Let $\theta_j = 2\pi j / M$, where $j = 1, 2, \dots, M$. The discrete coefficients are calculated as:
\begin{equation}
    f(\mathcal{L}) \approx \frac{1}{M} \sum_{j=1}^{M} \frac{f(z_j)}{z_j - \mathcal{L}} z_j, \quad \text{where } z_j = R e^{i\theta_j}
\end{equation}

In our specific initialization architecture, we use $M = 32$ complex quadrature points and a radius $R = 1/\Delta t$. This pre-computation is performed exactly once at $t=0$. The resulting robust coefficient tensors are stored in RAM, guaranteeing unconditionally stable and perfectly accurate integration of the Hall-MHD equations, fully exploiting the M4 processor's floating-point bandwidth without ever dividing by zero.

\section{Topological Invariants: Generalized Helicity Conservation}
\label{sec:helicity}

In order to rigorously train an Artificial Intelligence to predict topological collapse, we must possess an absolute, mathematically invariant ground truth against which we can measure the plasma's geometry. 

In standard Ideal MHD ($\eta = 0, d_i = 0$), this topological ground truth is the \textbf{Magnetic Helicity}:
\begin{equation}
    H = \int_V \mathbf{A} \cdot \mathbf{B} \, dV
\end{equation}
Magnetic helicity acts as a continuous measure of the knottedness, linking, and twisting of the magnetic field lines. Under Ideal MHD, the frozen-in flux theorem dictates that magnetic field lines move perfectly with the bulk fluid. Consequently, the topology cannot change, and $\frac{dH}{dt} = 0$. 

\subsection{The Hall-MHD Topological Symmetry Break}
However, our engine simulates \textit{Hall-MHD}. The inclusion of the Hall term ($\frac{d_i}{\rho} \mathbf{J} \times \mathbf{B}$) in Equation (8) decouples the electrons from the heavy ions at length scales smaller than $d_i$. Because the magnetic field is frozen into the \textit{electron} fluid rather than the bulk ion fluid, the macroscopic magnetic field lines are allowed to slip through the bulk plasma. 

Mathematically, this means standard magnetic helicity $H$ is \textbf{not} conserved in Hall-MHD. If we were to benchmark our Neural ODE's predictions against the decay of $H$, we would be measuring a mathematical artifact of the two-fluid separation, not true topological reconnection. 

\subsection{The Generalized Canonical Formulation}
To recover a conserved topological invariant, we must shift our perspective from pure electromagnetism to \textit{Hamiltonian mechanics}. In a two-fluid plasma, the true conserved quantities emerge when we consider the canonical momentum of the fluid species. 

Following Abdelhamid et al. (2015), we define a \textbf{Generalized Vector Potential} $\mathbf{A}^*$ that linearly combines the electromagnetic vector potential with the mechanical momentum of the ion fluid:
\begin{equation}
    \mathbf{A}^* = \mathbf{A} + d_i \mathbf{v}
\end{equation}
Taking the curl of this generalized potential yields the \textbf{Generalized Magnetic Field} $\mathbf{B}^*$:
\begin{equation}
    \mathbf{B}^* = \nabla \times \mathbf{A}^* = \mathbf{B} + d_i \boldsymbol{\omega}
\end{equation}
where $\boldsymbol{\omega} = \nabla \times \mathbf{v}$ is the macroscopic fluid vorticity. 

We can now define the \textbf{Generalized Helicity}, $H_{gen}$, which measures the interlocked topology of both the magnetic flux tubes and the fluid vortex tubes:
\begin{equation}
    H_{gen} = \int_V \mathbf{A}^* \cdot \mathbf{B}^* \, dV = \int_V (\mathbf{A} + d_i \mathbf{v}) \cdot (\mathbf{B} + d_i \boldsymbol{\omega}) \, dV
\end{equation}

\subsection{The Conservation Proof}
To prove that $H_{gen}$ is an absolute invariant of the collisionless ($\eta \to 0$) Hall-MHD system, we must show that its time derivative is strictly zero.

\textbf{The Non-Trivial Step:} Instead of expanding the full Hall-MHD induction and Navier-Stokes momentum equations (which results in a chaotic expansion of cross-products), we utilize a profound symmetry of the Hamiltonian formulation. In ideal Hall-MHD, the generalized vector potential $\mathbf{A}^*$ obeys an evolution equation perfectly analogous to ideal single-fluid MHD:
\begin{equation}
    \frac{\partial \mathbf{A}^*}{\partial t} = \mathbf{v} \times \mathbf{B}^* - \nabla \psi^*
\end{equation}
where $\psi^*$ is a modified scalar potential encompassing the pressure gradients and kinetic energy terms ($\frac{1}{2}|\mathbf{v}|^2$). 

Taking the curl of Equation (39), we obtain the evolution of the generalized magnetic field:
\begin{equation}
    \frac{\partial \mathbf{B}^*}{\partial t} = \nabla \times (\mathbf{v} \times \mathbf{B}^*)
\end{equation}

We now take the time derivative of the Generalized Helicity integral (Eq. 38). Using the product rule inside the integral gives:
\begin{equation}
    \frac{d H_{gen}}{dt} = \int_V \left( \frac{\partial \mathbf{A}^*}{\partial t} \cdot \mathbf{B}^* + \mathbf{A}^* \cdot \frac{\partial \mathbf{B}^*}{\partial t} \right) dV
\end{equation}

We evaluate the two terms inside the integrand separately. First, substitute Equation (39) into the first term:
\begin{equation}
    \frac{\partial \mathbf{A}^*}{\partial t} \cdot \mathbf{B}^* = (\mathbf{v} \times \mathbf{B}^*) \cdot \mathbf{B}^* - (\nabla \psi^*) \cdot \mathbf{B}^*
\end{equation}
Because the cross product $\mathbf{v} \times \mathbf{B}^*$ is orthogonal to $\mathbf{B}^*$, the dot product $(\mathbf{v} \times \mathbf{B}^*) \cdot \mathbf{B}^* \equiv 0$. Furthermore, using the vector calculus identity $\nabla \cdot (\phi \mathbf{V}) = \nabla \phi \cdot \mathbf{V} + \phi (\nabla \cdot \mathbf{V})$ and knowing $\nabla \cdot \mathbf{B}^* = 0$, we rewrite the remaining term:
\begin{equation}
    \frac{\partial \mathbf{A}^*}{\partial t} \cdot \mathbf{B}^* = - \nabla \cdot (\psi^* \mathbf{B}^*)
\end{equation}

Now we evaluate the second term of Equation (41). Substituting Equation (40) yields $\mathbf{A}^* \cdot \nabla \times (\mathbf{v} \times \mathbf{B}^*)$. We apply the vector identity $\nabla \cdot (\mathbf{X} \times \mathbf{Y}) = \mathbf{Y} \cdot (\nabla \times \mathbf{X}) - \mathbf{X} \cdot (\nabla \times \mathbf{Y})$. Letting $\mathbf{X} = \mathbf{A}^*$ and $\mathbf{Y} = \mathbf{v} \times \mathbf{B}^*$:
\begin{equation}
    \mathbf{A}^* \cdot \nabla \times (\mathbf{v} \times \mathbf{B}^*) = (\mathbf{v} \times \mathbf{B}^*) \cdot (\nabla \times \mathbf{A}^*) - \nabla \cdot \Big[ \mathbf{A}^* \times (\mathbf{v} \times \mathbf{B}^*) \Big]
\end{equation}
Recognizing that $\nabla \times \mathbf{A}^* = \mathbf{B}^*$, the first term on the right becomes $(\mathbf{v} \times \mathbf{B}^*) \cdot \mathbf{B}^*$, which is again exactly zero. Thus:
\begin{equation}
    \mathbf{A}^* \cdot \frac{\partial \mathbf{B}^*}{\partial t} = - \nabla \cdot \Big[ \mathbf{A}^* \times (\mathbf{v} \times \mathbf{B}^*) \Big]
\end{equation}

Substituting Equations (43) and (45) back into the global integral (Eq. 41) yields:
\begin{equation}
    \frac{d H_{gen}}{dt} = \int_V -\nabla \cdot \left( \psi^* \mathbf{B}^* + \mathbf{A}^* \times (\mathbf{v} \times \mathbf{B}^*) \right) dV
\end{equation}

\subsection{Causal Conclusion}
Equation (46) reveals that the time derivative of Generalized Helicity is exactly the volume integral of a pure divergence. By applying \textbf{Gauss's Divergence Theorem}, we convert this volume integral into a surface integral over the boundary $\partial V$:
\begin{equation}
    \frac{d H_{gen}}{dt} = \oint_{\partial V} - \left( \psi^* \mathbf{B}^* + \mathbf{A}^* \times (\mathbf{v} \times \mathbf{B}^*) \right) \cdot d\mathbf{S} = 0
\end{equation}
Because our C++ physics engine operates in a triply periodic Fourier domain ($T^3$), there are no physical boundaries; the fluxes exiting one face of the cube re-enter the opposite face identically, causing the surface integral to evaluate to exactly zero. 

Therefore, Generalized Helicity is rigorously conserved. Any decay in $H_{gen}$ observed during our numerical simulation is mathematically guaranteed to be the direct, causal consequence of the explicitly added hyper-resistivity ($\eta_4 \nabla^4 \mathbf{A}$) breaking the flux tubes. This gives us the ultimate topological ground truth: when the Artificial Intelligence predicts a singularity, it is predicting a genuine, physically irreversible thermodynamic event.

\newpage
\begin{thebibliography}{99}

\bibitem[Biskamp(2000)]{biskamp2000}
Biskamp, D. (2000). \textit{Magnetic Reconnection in Plasmas}. Cambridge University Press.

\bibitem[Peyret(2002)]{peyret2002}
Peyret, R. (2002). \textit{Spectral Methods for Incompressible Viscous Flow}. Springer.

\bibitem[Moffatt(1969)]{moffatt1969}
Moffatt, H. K. (1969). The degree of knottedness of a tangled magnetic field. \textit{Journal of Fluid Mechanics}, 35(1), 117-129.

\bibitem[Kassam \& Trefethen(2005)]{kassam2005}
Kassam, A. K., \& Trefethen, L. N. (2005). Fourth-order time-stepping for stiff PDEs. \textit{SIAM Journal on Scientific Computing}, 26(4), 1214-1233.

\bibitem[Abdelhamid et al.(2015)]{abdelhamid2015}
Abdelhamid, H. M., Kawazura, Y., \& Yoshida, Z. (2015). Hamiltonian formalism of extended magnetohydrodynamics. \textit{Journal of Physics A: Mathematical and Theoretical}, 48(23), 235502.

\end{thebibliography}

\end{document}